\section{Scope, Method, and Qualifications}

\paragraph{Reader and purpose.} This is spiritual reading for Catholics at the close of the liturgical year, the third of four volumes on the last things. It sets out the Church's faith concerning hell---its reality, its pains, its eternity, and the risen body's share in it---through Scripture, the councils, the liturgy, and the Fathers, Doctors and saints, and turns that consideration toward contrition, confession, perseverance and prayer. Death, judgment and heaven, purgatory and prayer for the dead, and the particular and general judgment are treated in the companion volumes and appear here only where the subject requires.

\paragraph{Grades.} Each teaching is stated in the body at its degree in a single clause; the appendix \emph{Teaching and Its Grades} gives the grade of each, using six terms: defined; Church teaching; common teaching of the Fathers and Doctors; theological opinion; pious tradition; approved private revelation. Where the Church leaves a question open, the appendix names the positions held.

\paragraph{Sources and translations.} Scripture is quoted from the Douay--Rheims (Challoner revision) in its Vulgate numbering, so that the refrain of Mark 9 stands at verses 43, 45 and 47. The Fathers are quoted from the nineteenth-century English translations of the \emph{Ante-Nicene} and \emph{Nicene and Post-Nicene Fathers}, the Dods translation of the \emph{City of God}, the Library of the Fathers translation of Gregory's \emph{Morals}, and the 1608 English of Gregory's \emph{Dialogues} as re-edited in 1911, whose chapter numbers are used. The \emph{Summa} and its Supplement are quoted in the translation of the Fathers of the English Dominican Province. Ignatius is quoted in Mullan's translation of 1914, which prints no paragraph numbers, and is cited by week and exercise. Teresa is quoted in David Lewis's translation; Alphonsus in the Dublin translation of 1843; Francis de Sales in the Rivingtons translation of 1876; Bernard in George Lewis's translation of 1908; Catherine of Siena in Algar Thorold's abridged edition of 1907; the \emph{Imitation} in William Benham's translation. Conciliar and papal Latin is quoted from the Denzinger numbering; Trent in English from Waterworth. Liturgical Latin is quoted from the \emph{Rituale Romanum} of 1872, the 1962 Missal, and hand missals that print the Missal's Latin; liturgical English from F. X. Lasance's \emph{New Roman Missal} (1945), the \emph{Catechism}'s citation of the Roman Canon, and W. J. Irons's \emph{Dies irae}. Where no public-domain English was available (Gregory's \emph{Homilies on the Gospels}, the Office of the Dead), the Latin is given and its sense is paraphrased in the surrounding prose, not offered as a translation. The closing prayer is reproduced from its identified witness without alteration.

\paragraph{Rights.} Texts of the Holy See (the \emph{Catechism}, the 1999 audience of John Paul II, Paul VI's \emph{Credo}, and the 2000 publication of the Fatima secret, including Sister Lúcia's memoir in its official English) remain the property of their publishers and are quoted briefly with attribution. Scripture, liturgical texts, received prayers, and third-party translations are not relicensed by this publication.

\paragraph{Limits.} The doctrinal history of the eternity of punishment, the debates of the schools, and modern controversies lie outside its scope, and no judgment is made of the state of any person beyond what Scripture, the liturgy and the Fathers say of Judas. It is a work of study and devotion; it is not a catechism, does not replace the teaching of the Church's pastors, and carries no ecclesiastical approval.
